\documentclass[11pt,a4paper]{article}
\usepackage[margin=25mm]{geometry}
\usepackage[T1]{fontenc}
\usepackage[utf8]{inputenc}
\usepackage{lmodern,amsmath,amssymb,enumitem,xurl,xcolor,needspace}
\usepackage[colorlinks=true,linkcolor=black,urlcolor=blue!55!black]{hyperref}
\ifsolutions
\hypersetup{pdfauthor={Luca Scharr},pdftitle={Tutorial 11 - From Images to Depth - Worked Solutions}}
\else
\hypersetup{pdfauthor={Luca Scharr},pdftitle={Tutorial 11 - From Images to Depth}}
\fi
\setlength{\parindent}{0pt}
\setlength{\parskip}{0.45em}
\setlist{leftmargin=*,itemsep=0.15em,topsep=0.15em}
\emergencystretch=2em
\newcommand{\laburl}{https://sosekie.github.io/Unibe-CVG-CV-Interactive-Labs/tutorial-11/}
\newcommand{\lablink}[2]{\href{\laburl\##1}{#2}}
\begin{document}
\begin{center}
{\Large\bfseries Tutorial 11 - From Images to Depth\ifsolutions\ - Worked Solutions\fi}\\[2mm]
Luca Scharr\\
Computer Vision 2026\\
University of Bern\\[1mm]
Interactive lab: \href{\laburl}{Open the monocular depth tutorial}
\end{center}
\ifsolutions
Each question is repeated before its worked answers.
\else
\section*{Learning objectives}
\begin{itemize}[nosep]
\item Interpret monocular predictions in light of single-view ambiguity.
\item Explain the training pipeline of Depth Anything V2 and assess evidence from its ablations.
\item Relate Depth Pro's multiscale architecture and focal-length estimate to its output.
\item Explain MoGe-2's point-map representation, scale prediction and training-data refinement.
\end{itemize}
Explain your reasoning and state your assumptions.
Use the interactive tutorial to inspect examples and check your predictions.
\fi
\needspace{10\baselineskip}
\section*{Part 1. Model playground}
\subsection*{1. What do depth maps tell us?}
\noindent\textbf{Interactive experiment.}
Open \lablink{play}{Model playground}, choose an indoor photograph and predict where depth will be difficult to estimate.
Reveal the predictions and compare them with the sensor reference.
The Depth Anything V2 output is a relative inverse-depth score; the MoGe-2 output is a metric-depth estimate.
\par\noindent
\begin{minipage}{\linewidth}
\subsubsection*{1.1}
A calibrated pinhole camera observes a point $(X,Y,Z)$, whose horizontal image coordinate is $u=fX/Z$.
\begin{enumerate}[label=(\alph*),nosep]
\item What happens to $u$ if the 3D coordinates are multiplied by a positive constant $k$?
Why does this matter for metric depth from a single image?
\item How can a trained model still produce a metric estimate, and why can that estimate be incorrect?
\end{enumerate}
\ifsolutions
\noindent\textbf{Solution (a).}
The projection is unchanged because $u'=f(kX)/(kZ)=fX/Z$.
A single image does not by itself determine distance along the camera ray or the absolute scale of the scene.

\noindent\textbf{Solution (b).}
Training data can provide priors and metric supervision that help the model choose a plausible scale.
Those learned regularities may not hold for the input scene, so a metric prediction is an estimate rather than a direct measurement.
\fi
\end{minipage}
\par\medskip
\par\noindent
\begin{minipage}{\linewidth}
\subsubsection*{1.2}
A relative model predicts inverse-depth scores $s_A=2$ and $s_B=5$ for two visible points.
Larger scores indicate nearer points.
\begin{enumerate}[label=(\alph*),nosep]
\item Which point is predicted nearer?
Is their distance in metres determined by these scores?
\item The visualisations independently normalise relative scores for display.
Why can equal colours in two maps not establish equal metric depth?
What additional evidence would help evaluate the prediction?
\end{enumerate}
\ifsolutions
\noindent\textbf{Solution (a).}
Point B is predicted nearer because its inverse-depth score is larger.
The scores do not provide a calibrated distance in metres.

\noindent\textbf{Solution (b).}
A display colour depends on the normalisation and output convention, so the same shade can represent different values in different maps.
Independent calibrated depth measurements or additional calibrated views can help assess the prediction.
Sensor references can also be missing or unreliable, particularly on reflective surfaces.
\fi
\end{minipage}
\par\medskip
\needspace{10\baselineskip}
\section*{Part 2. Depth Anything V2}
\subsection*{2. Learning from synthetic and real images}
\noindent\textbf{Interactive experiment.}
Open \lablink{data}{Depth Anything V2} and follow the four training stages.
The pipeline illustration and the displayed student prediction are not live training.
\par\noindent
\begin{minipage}{\linewidth}
\subsubsection*{2.1}
The method uses labelled synthetic images, a large teacher and unlabelled real images.
\begin{enumerate}[label=(\alph*),nosep]
\item Describe how the teacher and student obtain their respective training targets.
Why can the real photographs be used without sensor depth?
\item What does the real-image stage address, and what limitation can the generated targets introduce?
\end{enumerate}
\ifsolutions
\noindent\textbf{Solution (a).}
The teacher is trained on synthetic RGB images paired with depth from rendered geometry.
It predicts depth for real photographs, producing pseudo-labels that serve as targets for a student model.
The real photographs do not require measured depth for this stage.

\noindent\textbf{Solution (b).}
Real images provide diverse appearances and help address the gap between synthetic and real image distributions.
Pseudo-labels can contain systematic or local teacher errors, which may be transferred to the student.
They are not independent ground-truth measurements.
\fi
\end{minipage}
\par\medskip
\par\noindent
\begin{minipage}{\linewidth}
\subsubsection*{2.2}
The tutorial reproduces the paper's DA-2K results for a ViT-S student: 89.8\% accuracy with synthetic labels only and 95.3\% with pseudo-labelled real images only.
DA-2K assesses relative depth ordering on annotated point pairs.
\begin{enumerate}[label=(\alph*),nosep]
\item Calculate the difference in accuracy in percentage points.
What additional training did the teacher require in the second condition?
\item Does this comparison establish that every pseudo-label is correct or that the model produces metric depth?
Explain.
\end{enumerate}
\ifsolutions
\noindent\textbf{Solution (a).}
The reported difference is $95.3-89.8=5.5$ percentage points.
The teacher producing the pseudo-labels was trained on labelled synthetic images.

\noindent\textbf{Solution (b).}
No.
The score aggregates relative-ordering results under the benchmark protocol, so it does not validate each pseudo-label.
The relative model's scores are not calibrated metre values.
Other datasets and criteria are needed to assess different failure modes and claims.
\fi
\end{minipage}
\par\medskip
\needspace{10\baselineskip}
\section*{Part 3. Depth Pro}
\subsection*{3. Context, detail and metric depth}
\noindent\textbf{Interactive experiment.}
Open \lablink{detail}{Depth Pro} and select a location on a fine structure in the bird-cage or cat image.
Compare the aligned crops and follow the multiscale pipeline.
The input and outputs are visualisations from the Depth Pro paper, while the feature maps in the pipeline diagram are schematic.
\par\noindent
\begin{minipage}{\linewidth}
\subsubsection*{3.1}
Depth Pro processes three image scales using 1, 9 and 25 image patches.
The two finer scales use overlapping patches, and the shared patch encoder produces a $24\times24$ grid of feature locations per patch.
\begin{enumerate}[label=(\alph*),nosep]
\item How many patches are processed in total, and how many feature locations do these grids contain before merging?
\item Why can local patches and larger-scale context be useful together?
Why should overlapping regions not be counted as independent image evidence?
\end{enumerate}
\ifsolutions
\noindent\textbf{Solution (a).}
There are $1+9+25=35$ patches and $35\cdot24\cdot24=20160$ feature locations before merging.
This count is for the stated patch grids and excludes additional intermediate features and the separate image encoder.

\noindent\textbf{Solution (b).}
Fine patches retain local spatial detail, while a larger field of view can provide cues about scene layout and object identity.
Overlapping patches contain some of the same pixels, so their features are not independent observations.
The diagram illustrates the architecture but does not show the model's actual internal activations.
\fi
\end{minipage}
\par\medskip
\par\noindent
\begin{minipage}{\linewidth}
\subsubsection*{3.2}
Depth Pro computes metric depth using $D=f/(wC)$, where $f$ is focal length in pixels, $w$ is image width in pixels and $C$ is predicted canonical inverse depth.
For this question, hold $w=1000$ and $C=0.5$ fixed.
\begin{enumerate}[label=(\alph*),nosep]
\item Compute $D$ for $f=1000$ pixels and for $f=750$ pixels.
\item What does a focal-length estimate contribute, and why does it not remove the fundamental single-image ambiguity?
\end{enumerate}
\ifsolutions
\noindent\textbf{Solution (a).}
For $f=1000$, $D=1000/(1000\cdot0.5)=2$ m.
For $f=750$, $D=750/(1000\cdot0.5)=1.5$ m under the stated fixed-output assumption.

\noindent\textbf{Solution (b).}
The focal length supplies a camera-dependent factor in the metric conversion and can be estimated when camera metadata are unavailable.
Predicted focal length and inverse depth can both be wrong.
They depend on learned information and do not make the scene uniquely determined by one image.
Changing the input camera may also change the network output, so holding $C$ fixed is only an algebraic comparison.
\fi
\end{minipage}
\par\medskip
\par\noindent
\begin{minipage}{\linewidth}
\subsubsection*{3.3}
The tutorial reports Hypersim boundary F1 scores of 0.465 and 0.461 for two second-stage training-loss configurations in the Depth Pro paper.
Higher boundary F1 is better.
\begin{enumerate}[label=(\alph*),nosep]
\item Which setting has the higher reported score, and what is the absolute difference?
\item Can this loss ablation establish that the multiscale patch architecture caused the difference?
Explain.
\end{enumerate}
\ifsolutions
\noindent\textbf{Solution (a).}
The configuration with first- and second-order derivative losses scores 0.465, compared with 0.461 without derivative losses.
The difference is 0.004 F1.

\noindent\textbf{Solution (b).}
No.
The comparison tests the stated training-loss configurations, not the effect of removing the multiscale architecture.
An architectural ablation would be needed to isolate that design choice.
The paper reports further datasets and metrics, so this result does not describe all performance.
\fi
\end{minipage}
\par\medskip
\needspace{10\baselineskip}
\section*{Part 4. MoGe-2}
\subsection*{4. Point maps, scale and supervision}
\noindent\textbf{Interactive experiment.}
Open \lablink{geometry}{MoGe-2 in 3D} and rotate the predicted point map.
Compare it with the reference back-projection, then inspect the scale and annotation-refinement examples.
\par\noindent
\begin{minipage}{\linewidth}
\subsubsection*{4.1}
A depth map gives an axial depth $Z$ at pixel $(u,v)$ for a camera with intrinsics $(f_x,f_y,c_x,c_y)$.
A point map supplies a 3D point at each represented image location.
\begin{enumerate}[label=(\alph*),nosep]
\item Write the 3D camera-coordinate point obtained by back-projecting the pixel.
\item Why does rotating a plausible point map not establish the geometry behind an occluder?
Why should differences between reference and prediction not automatically be attributed to depth error alone?
\end{enumerate}
\ifsolutions
\noindent\textbf{Solution (a).}
The back-projected point is
\[
(X,Y,Z)=\left(\frac{(u-c_x)Z}{f_x},\frac{(v-c_y)Z}{f_y},Z\right).
\]
This assumes axial depth and the stated pinhole camera convention.

\noindent\textbf{Solution (b).}
The map represents points associated with visible image locations and does not independently observe hidden surfaces.
The displayed reference is back-projected with fixed NYU camera intrinsics, while the MoGe-2 point map is predicted by the model.
Differences can reflect depth, camera assumptions and other reconstruction errors, as well as imperfections in the sensor reference.
\fi
\end{minipage}
\par\medskip
\par\noindent
\begin{minipage}{\linewidth}
\subsubsection*{4.2}
MoGe-2 separates an affine-invariant point-map prediction from a global scale prediction.
For a simplified illustration, a relative point is $(0.5,0.3,2.0)$ and a positive scale is $s=1.5$.
\begin{enumerate}[label=(\alph*),nosep]
\item Multiply the point by $s$.
What happens to its projection when it is scaled about the camera centre?
\item Why might separate relative-geometry and scale branches be useful?
What part of the full representation does this scaling-only illustration omit?
\end{enumerate}
\ifsolutions
\noindent\textbf{Solution (a).}
The scaled point is $(0.75,0.45,3.0)$.
Its projection remains unchanged because the ratios $X/Z$ and $Y/Z$ remain unchanged.

\noindent\textbf{Solution (b).}
Separate branches allow the model to learn relative geometry while estimating global metric scale using global image features.
The separation is intended to reduce interference between these objectives, but neither prediction is guaranteed to be correct.
The affine-invariant representation also involves a global alignment shift, which the scaling-only diagram omits.
\fi
\end{minipage}
\par\medskip
\par\noindent
\begin{minipage}{\linewidth}
\subsubsection*{4.3}
Real training depth can be incomplete or misaligned near object boundaries.
MoGe-2 filters suspect measurements and completes missing regions using guidance from a model trained on synthetic data.
The tutorial reports metric point errors of 9.01\% and 9.20\%, and boundary F1 scores of 10.3 and 12.5, for raw and improved real-data configurations respectively.
\begin{enumerate}[label=(\alph*),nosep]
\item How can filtering and completion help with the example annotation?
Why is the completed label not automatically independent ground truth?
\item Compare the reported metric point error and boundary F1 for the two configurations.
What can and cannot be concluded from these aggregate values?
\end{enumerate}
\ifsolutions
\noindent\textbf{Solution (a).}
Filtering can remove mismatched depth values, and completion can fill gaps using a synthetic-trained model's local geometry guidance.
The resulting annotation depends on a prediction and can itself contain errors.
The true profile in the schematic is assumed for illustration and is not generally available for real data.

\noindent\textbf{Solution (b).}
The improved-data configuration has a slightly higher metric point error, 9.20\% versus 9.01\%, and a higher boundary F1, 12.5 versus 10.3.
Under the paper's evaluation, this suggests a boundary-quality gain with similar aggregate metric error.
It does not establish that every completed pixel is correct or that the result transfers to every scene.
The two metrics are aggregated over different sets of datasets.
\fi
\end{minipage}
\par\medskip
\section*{Papers}
\begin{itemize}[nosep]
\item Yang et al., \emph{Depth Anything V2}, NeurIPS 2024.
\href{https://arxiv.org/abs/2406.09414}{Paper}
\item Bochkovskii et al., \emph{Depth Pro: Sharp Monocular Metric Depth in Less Than a Second}, ICLR 2025.
\href{https://arxiv.org/abs/2410.02073}{Paper}
\item Wang et al., \emph{MoGe-2: Accurate Monocular Geometry with Metric Scale and Sharp Details}, NeurIPS 2025.
\href{https://arxiv.org/abs/2507.02546}{Paper}
\end{itemize}
\end{document}
